\documentclass[UTF8,12pt,a4paper,oneside,fontset=none]{ctexart}

\usepackage[a4paper,top=2.35cm,bottom=2.35cm,left=2.45cm,right=2.45cm,headheight=15pt]{geometry}
\usepackage{fontspec}
\usepackage{amsmath,amssymb,mathtools,bm}
\usepackage{unicode-math}
\usepackage{booktabs,longtable,tabularx,array,multirow}
\usepackage[table]{xcolor}
\usepackage{graphicx,float,subcaption}
\usepackage{siunitx}
\usepackage{enumitem}
\usepackage{fancyhdr}
\usepackage{caption}
\usepackage{hyperref}
\usepackage[nameinlink,noabbrev]{cleveref}
\usepackage[most]{tcolorbox}
\usepackage{listings}
\usepackage{seqsplit}
\usepackage{microtype}

\setmainfont{Times New Roman}
\setsansfont{Arial}
\setmonofont{Consolas}
\setmathfont{Cambria Math}
\setCJKmainfont[BoldFont=SimHei,ItalicFont=KaiTi]{SimSun}
\setCJKsansfont[BoldFont={Microsoft YaHei Bold}]{Microsoft YaHei}
\setCJKmonofont{FangSong}

\definecolor{DeepBlue}{HTML}{17324D}
\definecolor{PandaTeal}{HTML}{0B7A75}
\definecolor{WarmGold}{HTML}{C58A19}
\definecolor{SoftBlue}{HTML}{EEF5F8}
\definecolor{SoftTeal}{HTML}{ECF7F5}
\definecolor{SoftGold}{HTML}{FFF7E7}
\definecolor{SoftRed}{HTML}{FFF0F0}
\definecolor{MidGray}{HTML}{6B7280}
\definecolor{LineGray}{HTML}{D7DEE5}

\hypersetup{
  colorlinks=true,
  linkcolor=DeepBlue,
  urlcolor=PandaTeal,
  citecolor=PandaTeal,
  bookmarksnumbered=true,
  pdfauthor={范子恒},
  pdftitle={基于 PandaX-4T 标定数据的独立能量重建与分辨率改善研究}
}

\crefname{equation}{式}{式}
\Crefname{equation}{式}{式}
\crefname{table}{表}{表}
\Crefname{table}{表}{表}
\crefname{figure}{图}{图}
\Crefname{figure}{图}{图}

\ctexset{
  section={format=\Large\bfseries\color{DeepBlue},beforeskip=1.5ex,afterskip=1.0ex},
  subsection={format=\large\bfseries\color{PandaTeal},beforeskip=1.2ex,afterskip=0.7ex},
  subsubsection={format=\normalsize\bfseries\color{DeepBlue}}
}

\pagestyle{fancy}
\fancyhf{}
\fancyhead[L]{\small PandaX-4T 独立能量重建研究记录}
\fancyhead[R]{\small\nouppercase{\leftmark}}
\fancyfoot[C]{\small\color{MidGray}\thepage}
\renewcommand{\headrulewidth}{0.4pt}
\renewcommand{\footrulewidth}{0pt}

\captionsetup{font=small,labelfont=bf,labelsep=quad}
\graphicspath{{figures/}}
\setlength{\parindent}{2em}
\setlength{\parskip}{0.25em}
\renewcommand{\arraystretch}{1.18}
\setcounter{tocdepth}{2}
\numberwithin{equation}{section}

\sisetup{
  round-mode=places,
  round-precision=3,
  group-separator={,},
  group-minimum-digits=4
}

\DeclareMathOperator{\median}{median}
\DeclareMathOperator{\quantile}{Quantile}
\DeclareMathOperator*{\argminop}{arg\,min}
\newcommand{\Sone}{\mathrm{S1}}
\newcommand{\Stwo}{\mathrm{S2}}
\newcommand{\StwoB}{\mathrm{S2}_{\mathrm{B}}}
\newcommand{\Rsixtyeight}{R_{68}}
\newcommand{\Rninety}{R_{90}}
\newcommand{\OOF}{\mathrm{OOF}}
\newcommand{\Erel}{\hat e}
\newcommand{\dd}{\mathrm{d}}

\newcolumntype{Y}{>{\raggedright\arraybackslash}X}
\newcolumntype{C}[1]{>{\centering\arraybackslash}p{#1}}
\newcolumntype{R}[1]{>{\raggedleft\arraybackslash}p{#1}}

\newtcolorbox{questionbox}[1][]{
  enhanced,breakable,
  colback=SoftBlue,colframe=DeepBlue,
  boxrule=0.7pt,arc=2mm,
  title={本章要回答的问题},fonttitle=\bfseries,
  #1
}
\newtcolorbox{conclusionbox}[1][]{
  enhanced,breakable,
  colback=SoftTeal,colframe=PandaTeal,
  boxrule=0.7pt,arc=2mm,
  title={本章得到的结论},fonttitle=\bfseries,
  #1
}
\newtcolorbox{warningbox}[1][]{
  enhanced,breakable,
  colback=SoftRed,colframe=red!65!black,
  boxrule=0.7pt,arc=2mm,
  title={必须注意},fonttitle=\bfseries,
  #1
}
\newtcolorbox{resultbox}[1][]{
  enhanced,breakable,
  colback=SoftGold,colframe=WarmGold,
  boxrule=0.9pt,arc=2mm,
  title={核心结果},fonttitle=\bfseries,
  #1
}

\lstset{
  basicstyle=\ttfamily\footnotesize,
  breaklines=true,
  frame=single,
  rulecolor=\color{LineGray},
  backgroundcolor=\color{SoftBlue},
  columns=fullflexible,
  showstringspaces=false,
  upquote=true
}

\begin{document}

\begin{titlepage}
  \centering
  \vspace*{1.5cm}
  {\color{DeepBlue}\rule{\textwidth}{1.2pt}}\\[1.2cm]
  {\zihao{1}\bfseries\color{DeepBlue}
  基于 PandaX-4T 标定数据的\\[0.35cm]
  独立能量重建与分辨率改善研究\par}
  \vspace{0.8cm}
  {\zihao{3}\color{PandaTeal}
  仅使用基础观测量的响应校正、交叉验证与峰形审计\par}
  \vspace{1.2cm}
  {\color{DeepBlue}\rule{0.72\textwidth}{0.6pt}}\\[1.2cm]

  \begin{tcolorbox}[
    width=0.82\textwidth,
    colback=SoftBlue,colframe=DeepBlue,
    arc=2mm,boxrule=0.7pt]
    \centering
    \begin{tabular}{rl}
      项目成员：& 范子恒 \\
      数据集：& PandaX-4T run 10972 \\
      报告性质：& 可复核的研究过程与阶段性结果记录 \\
      报告日期：& 2026 年 7 月 15 日 \\
      版本：& v1.0（独立重建阶段总结）
    \end{tabular}
  \end{tcolorbox}

  \vfill
  \begin{minipage}{0.86\textwidth}
    \small\color{MidGray}
    本报告完整记录从原始 TXT 文件审计、历史方法复核、独立协议冻结、
    负结果分析、物理可解释主方案、组合响应候选方案、文件级交叉验证、
    bootstrap 到峰形拟合优度审计的全部关键步骤。报告中的百分比宽度均明确标注定义，
    未通过拟合优度检验的 Gaussian 核心宽度不作为正式成果。
  \end{minipage}
  \vspace{0.8cm}
  {\color{DeepBlue}\rule{\textwidth}{1.2pt}}
\end{titlepage}

\pagenumbering{Roman}

\section*{摘要}
\addcontentsline{toc}{section}{摘要}

本研究的目标，是在不依赖以前研究者重建算法和已修正能量分支的条件下，
仅根据当前数据中的基础观测量，探索能否稳定改善 PandaX-4T MeV 能区样本的相对能量重建宽度。
输入为 run 10972 的一个 TXT 文件，共包含 \(4500\) 个事件和 \(188\) 个字段。
正式独立流程只使用原始 \(\Sone\)、底部 \(\StwoB\)、漂移时间、原始 S2 位置以及事件标识，
不应用语义未确认的 cut，不读取已有的 corrected、lifetime、desaturation 或 MCPAF 分支，
并保留全部 \(4500\) 个事件。

研究首先复核了早期派生分支分析，发现旧 Gaussian 加指数背景模型对拟合窗口高度敏感，
且正式参数 bootstrap 拟合优度检验给出 \(p=0.004975\)，因此旧的窄核心宽度被判定为不可靠。
随后冻结独立变量白名单，建立按文件分组的嵌套五折交叉验证。
物理可解释的双通道 A4 方案分别修正 \(\Sone\) 与 \(\StwoB\) 的漂移时间和径向响应，
在连续文件块测试中将全事件非参数宽度 \(\Rsixtyeight\) 从 \(7.893\%\) 降至 \(6.779\%\)，
将 \(\Rninety\) 从 \(14.783\%\) 降至 \(13.278\%\)。
更灵活的组合响应 challenger 将同一拆分下的 \(\Rsixtyeight\) 从 \(8.125\%\) 降至 \(5.973\%\)，
将 \(\Rninety\) 从 \(15.029\%\) 降至 \(11.826\%\)。两套方法的五个连续外层测试块均改善。

当前结果证明：仅使用基础变量的独立相对能量重建，在本 run 的未见文件块上存在稳定改善空间。
但是，输入文件名包含 \texttt{Egt2MeV}，说明数据可能经历未知能量阈值预选；同时目前只有一个 run，
没有绝对峰能、\(g_1\) 或 \(g_2\) 标定，也没有完全独立的盲测样本。
因此本报告不把上述数字表述为最终物理能量分辨率，而将 A4 冻结为主方案、challenger 冻结为待验证候选。

\begin{resultbox}[title={一句话成果}]
在不删除事件、不使用已有校正分支的条件下，物理可解释主方案在连续文件块 OOF 测试中实现
\(\Rsixtyeight:7.893\%\rightarrow6.779\%\)；更灵活候选实现
\(8.125\%\rightarrow5.973\%\)。这是一项当前 run 上的独立相对重建改善，
还不是跨 run、无选择偏差的最终能量分辨率测量。
\end{resultbox}

\clearpage
\tableofcontents
\clearpage
\pagenumbering{arabic}

\section{给初学者的阅读指南}

\begin{questionbox}
这项工作到底在解决什么问题？为什么“谱更窄”不一定就代表方法更好？本报告应该按照什么顺序阅读？
\end{questionbox}

\subsection{最短的理解路径}

一次粒子相互作用在液氙中沉积能量后，会产生闪烁光和电离电子。
PandaX-4T 把它们分别记录为 \(\Sone\) 和 \(\Stwo\)。理想情况下，相同能量的事件应重建到相同位置；
现实中，电子漂移损失、光收集效率、位置依赖和探测器随时间变化会使同一能量对应一段分布。
本工作的核心就是：只利用训练数据估计这些响应变化，在未参加训练的数据上做校正，
再检查分布是否真正变窄、尾部是否同时改善。

\begin{center}
\begin{tcolorbox}[width=0.92\textwidth,colback=white,colframe=LineGray,boxrule=0.7pt,arc=1mm]
\centering
粒子沉积能量
\(\longrightarrow\) 产生光子和电子
\(\longrightarrow\) 观测到 \(\Sone\) 与 \(\StwoB\)\\[0.35em]
\(\longrightarrow\) 学习漂移时间和空间响应
\(\longrightarrow\) 校正并组合两个通道
\(\longrightarrow\) 在未见文件上评价谱宽
\end{tcolorbox}
\end{center}

\subsection{七个必须先懂的词}

\begin{description}[style=nextline,leftmargin=2.7cm,labelwidth=2.3cm]
  \item[事件（event）] 一次候选物理相互作用，对应数据表中的一行。
  \item[run] 一段连续取数。本报告只有 run 10972，因此无法直接检验跨 run 稳定性。
  \item[file] run 内部的数据文件。交叉验证按 file 分组，避免同一文件的信息同时进入训练和测试。
  \item[响应校正] 学习信号中心如何随漂移时间或位置变化，再除去这种系统性变化。
  \item[分辨率] 重建量的相对宽度，不是分类任务中的“准确率”。
  \item[OOF] out-of-fold。每个事件的预测都来自没有见过该事件所属外折的模型。
  \item[盲测] 方法和参数冻结后，在此前完全未接触的数据上只评价一次。
\end{description}

\subsection{如何判断一项改善是否可信}

本报告要求一项结论同时回答四个问题：

\begin{enumerate}[label=\arabic*.]
  \item 改善是否出现在未参与训练的文件上？
  \item 覆盖中心约 68\% 事件的相对宽度 \(\Rsixtyeight\)，以及覆盖中心 90\% 事件、
  对长尾更敏感的 \(\Rninety\)，是否同时改善？
  \item 是否保留固定事件集合，而不是删除难重建事件来得到窄峰？
  \item 如果引用参数化的 \(\sigma/\mu\)，峰形模型是否通过拟合优度检验？
\end{enumerate}

\begin{conclusionbox}
读者可以先看第~\ref{sec:metrics}、\ref{sec:timeline}、\ref{sec:main-results}、\ref{sec:final}~节，
再回到方法章节理解细节。
本报告最重要的纪律是：训练集上的漂亮数字不算结果，未通过 GOF 的核心宽度也不算结果。
\end{conclusionbox}

\section{双相液氙 TPC 与能量重建背景}

\begin{questionbox}
什么是 \(\Sone\) 和 \(\Stwo\)？为什么把两者组合通常比只看一个通道更合理？漂移时间和位置为什么会使能量谱变宽？
\end{questionbox}

\subsection{能量沉积、光子与电子}

粒子在液氙中沉积能量 \(E_{\mathrm{dep}}\) 后，能量主要转化为激发和电离产生的可探测量子。
常用的概念关系为

\begin{equation}
N_q
=
\frac{E_{\mathrm{dep}}}{W}
=
N_{\gamma}+N_e,
\label{eq:quanta}
\end{equation}

其中 \(N_{\gamma}\) 是闪烁光子数，\(N_e\) 是逃逸复合并被漂移电场收集的电子数，
\(W\) 是平均产生一个量子所需的能量，液氙中常取约 \(\SI{13.7}{eV}\)\cite{pandax-bb2nu}。
这里的关系用于解释重建思路；本数据没有提供足够信息重新测量 \(W\)。

\subsection{S1：瞬发闪烁信号}

液氙中的瞬发闪烁光被上下光电倍增管阵列观测，形成 S1。理想化写作

\begin{equation}
\Sone=g_1N_{\gamma},
\label{eq:s1}
\end{equation}

其中 \(g_1\) 汇总光收集、量子效率和电子学响应。真实探测器中的 \(g_1\) 具有空间依赖，
因此相同 \(N_{\gamma}\) 在不同位置可能产生不同的 S1 面积。

\subsection{S2：漂移电子产生的二次电致发光}

电离电子在电场中向液面漂移，经液气界面抽取到气相后产生二次电致发光，形成 S2。
本研究使用底部阵列观测到的 S2，即

\begin{equation}
\StwoB=g_{2,\mathrm B}N_e,
\label{eq:s2}
\end{equation}

其中下标 \(\mathrm B\) 表示 bottom。选择底部 S2 的原因是当前任务关注基础变量，
并希望降低顶部阵列位置依赖与可能的高面积非线性影响；这并不意味着底部 S2 完全没有位置依赖。

\subsection{为什么要组合 S1 与 S2}

电子复合会使 S1 和 S2 近似反相关：更多电子复合通常增加闪烁、减少逃逸电子。
若 \(g_1\) 和 \(g_{2,\mathrm B}\) 已知，理想的绝对组合能量可写为

\begin{equation}
E_{\mathrm{comb}}
=
W\left(
\frac{\Sone}{g_1}
+
\frac{\StwoB}{g_{2,\mathrm B}}
\right).
\label{eq:combined-energy}
\end{equation}

本阶段没有独立测得 \(g_1\)、\(g_{2,\mathrm B}\) 和目标峰能量，因此实际得到的是无量纲相对能量代理量 \(\Erel\)，
不能直接标成 MeV。报告中所有宽度都相对于各训练折的鲁棒中心归一化。

\subsection{漂移时间和空间位置为什么重要}

由 S1 与 S2 的时间差得到漂移时间 \(t_{\mathrm d}\)，它近似对应深度：

\begin{equation}
z\simeq v_{\mathrm d}t_{\mathrm d},
\label{eq:z-drift}
\end{equation}

其中 \(v_{\mathrm d}\) 为电子漂移速度。本分析只把 \(t_{\mathrm d}\) 当作响应坐标，
没有用\cref{eq:z-drift}做绝对 \(z\) 标定。横向半径定义为

\begin{equation}
r^2=x^2+y^2.
\label{eq:r2}
\end{equation}

电子在漂移过程中会被杂质俘获，最简单的寿命模型为

\begin{equation}
\StwoB^{\mathrm{obs}}(t_{\mathrm d})
=
\StwoB^{0}
\exp\left(-\frac{t_{\mathrm d}}{\tau_e}\right),
\label{eq:lifetime}
\end{equation}

其中 \(\tau_e\) 是电子寿命。光收集、气液界面、场形和位置重建也会引入 \(r\)、\(x\)、\(y\) 依赖。

\begin{conclusionbox}
物理上应同时处理 S1 与 S2B 的响应，再组合两个通道。本阶段只做相对能量重建；
任何写成 MeV 的绝对结果都超出了当前数据能够支持的范围。
\end{conclusionbox}

\section{研究目标、独立性原则与分析边界}

\begin{questionbox}
“不依赖以前的工作”在本项目中如何落实？哪些变量允许使用，哪些变量明确禁止？什么结论从一开始就不在本阶段声称？
\end{questionbox}

\subsection{研究问题}

本项目冻结的研究问题是：

\begin{quote}
能否不读取以前研究者的重建代码和已修正信号分支，仅根据当前数据中的基础观测变量，
学习一种可以在未见文件上稳定减小相对能量谱宽的方法？
\end{quote}

这里的“独立”是算法输入和模型选择意义上的独立，不是说完全忽略探测器物理。
变量说明文档可用于理解单位和物理含义；已有校正量、已有能量估计及其参数不能进入训练或选择。

\subsection{允许字段白名单}

正式独立流程只读取以下八个字段：

\begin{equation}
\mathcal V_{\mathrm{allow}}
=
\left\{
\begin{array}{c}
\texttt{runNumber},\ \texttt{fileNumber},\ \texttt{eventNumber},\\
\texttt{qS1\_max},\ \texttt{qS2B\_max},\ \texttt{dt},\\
\texttt{xS2T\_max},\ \texttt{yS2T\_max}
\end{array}
\right\}.
\label{eq:whitelist}
\end{equation}

身份字段用于去重和分组；物理建模实际使用后五个观测量。
\texttt{dt} 原始单位为 ns，分析中除以 \(1000\) 转为 \(\mu\mathrm s\)；位置单位为 mm，信号面积单位为 p.e.

\subsection{开发阶段禁用的信息}

\begin{itemize}
  \item 名称中含 \texttt{C}、\texttt{cor}、\texttt{Belife}、\texttt{ub}、\texttt{stretch} 的已有校正分支；
  \item 名称中含 \texttt{des}、\texttt{MCPAF}、\texttt{Cs} 的去饱和或高级重建分支；
  \item 已有能量分支、以前研究者的算法代码和已调参数；
  \item 语义未确认的 cut flags；
  \item 目标峰的假定绝对能量，及任何使用测试集标签直接回归到常数的做法。
\end{itemize}

\subsection{预先规定的结论边界}

本阶段允许声称“当前 run 上的 OOF 相对宽度改善”，但不允许声称：

\begin{itemize}
  \item 已完成绝对 MeV 能标；
  \item 已得到 PandaX-4T 的最终能量分辨率；
  \item 已证明方法能跨 run、跨源位置或跨能量迁移；
  \item 已消除未知上游预选造成的选择偏差。
\end{itemize}

\begin{conclusionbox}
白名单和结论边界把“独立探索”变成可检查的技术规则，而不是一句口号。
后续 A4 与 challenger 均遵守\cref{eq:whitelist}，历史派生分支探索只作为被审计的背景记录。
\end{conclusionbox}

\section{数据获取与完整性审计}

\begin{questionbox}
输入文件是否完整？事件是否唯一？哪些字段存在缺失、无穷或常数问题？为什么没有直接使用数据里的 cut flags？
\end{questionbox}

\subsection{步骤 1：下载、落盘与文件指纹}

输入文件为

\begin{center}
\path{light_ana_run10972_finalSS_Egt2MeV_scalar.txt}
\end{center}

文件大小为 \(6{,}426{,}705\) bytes。下载后计算 SHA-256：

\begin{center}
\small\ttfamily
\seqsplit{d026e0e30bf6d6af44f3b04848c96815bbbc062cf13f7017e2422dc054cf9d0b}
\end{center}

该哈希与最终实验矩阵的 \path{basic_matrix_outputs/metadata.json} 一致，说明后续分析使用的是同一原始文件。
报告不保存、不展示任何账户密码或登录凭据。

\subsection{步骤 2：行列、事件身份与保留率审计}

\begin{table}[H]
\centering
\caption{原始数据完整性概览。事件键重复数为零，正式独立分析保留全部事件。}
\label{tab:data-audit}
\begin{tabularx}{0.88\textwidth}{Y R{3.5cm}}
\toprule
项目 & 结果 \\
\midrule
总行数 & 4501（含 1 行表头） \\
事件数 & 4500 \\
字段数 & 188 \\
唯一 run & 10972 \\
唯一 \texttt{fileNumber} & 985（范围 0--999） \\
重复 \((\texttt{run},\texttt{file},\texttt{event})\) & 0 \\
独立流程有限事件 & 4500 \\
独立流程保留率 & 100\% \\
\bottomrule
\end{tabularx}
\end{table}

正式白名单中的五个物理观测量在全部 \(4500\) 个事件上均有限。
此外，\texttt{qS1\_max} 与 \texttt{qS2B\_max} 全为正，因此没有必要为数值合法性删除事件。

\subsection{步骤 3：基础变量分布检查}

\begin{table}[H]
\centering
\caption{正式独立分析使用的基础观测量范围。}
\label{tab:variable-range}
\begin{tabularx}{0.9\textwidth}{l Y l}
\toprule
字段 & 范围或典型值 & 单位 \\
\midrule
\texttt{qS1\_max} & 4975.750--19834.627；中位数约 11123.3 & p.e. \\
\texttt{qS2B\_max} & 49352.891--820462.562；中位数约 424060 & p.e. \\
\texttt{dt}/1000 & 10.112--842.300；中位数约 417.8 & \(\mu\mathrm s\) \\
\texttt{xS2T\_max} & $-377.087$--379.639 & mm \\
\texttt{yS2T\_max} & $-411.257$--382.957 & mm \\
\bottomrule
\end{tabularx}
\end{table}

原始半径中位数约为 \(359.1\,\mathrm{mm}\)，四分位区间约为
\(345.7\)--\(368.7\,\mathrm{mm}\)，说明事件在横向上高度不均匀，符合外源样本可能集中在探测器边缘的直觉。
每个 file 的事件数中位数仅为 4，最少 1、最多 14；因此不能把单个 file 的峰中心当作可靠标签。

\begin{figure}[H]
\centering
\includegraphics[width=0.98\textwidth]{data_overview.png}
\caption{数据概览。S1--S2B 相关性、漂移时间分布、横向位置覆盖和每文件事件数共同表明：
样本规模有限且空间分布不均匀，响应模型必须保持低自由度，并按文件分组验证。}
\label{fig:data-overview}
\end{figure}

\subsection{步骤 4：188 个字段的缺失、无穷和常数审计}

九个字段全部为 NaN：

\begin{quote}\small
\path{xS2max_cdfTM}, \path{yS2max_cdfTM},
\path{xS2max_cdfTMs}, \path{yS2max_cdfTMs},
\path{lS2max_cdfTM}, \path{rms_cdfTM},
\path{xS2max_cdfPAF}, \path{yS2max_cdfPAF},
\path{lhfS2max_cdfPAF}。
\end{quote}

\path{qS1C_max} 及相关分支各含 2 个无穷值；若干 MCPAF 分支含 1--3 个无穷值。
因此早期派生分支共同比较只能保留 \(4497\) 个事件，而基础变量流程可以保留全部 \(4500\) 个事件。
另外，\path{xS2max_desCorCog_maxS2}、对应的 y 分支以及 firstS2 版本均为常数 0，不能提供位置信息。

\subsection{步骤 5：cut 字段语义审计}

数据中 14 个 cut 字段的 0/1 方向明显不统一。例如 \texttt{basicCut} 和 \texttt{S1PerPmtCut} 全为 1，
\texttt{fv\_extendCut}、\texttt{gas\_s2Cut} 和 \texttt{wallCut} 全为 0；
\texttt{S1PatternCut} 有 4472 个 1，\texttt{deadtime\_enhanceCut} 则有 4468 个 0。
在没有权威说明“0 或 1 哪个表示通过”的情况下，擅自叠加这些 cut 可能改变事件集合并人为缩窄谱宽。
因此独立流程不应用任何额外 cut。

\begin{warningbox}
文件名包含 \texttt{Egt2MeV}，说明样本很可能已经经过某种 \(E>2\,\mathrm{MeV}\) 预选。
当前不知道预选使用了哪一种能量估计，也不知道它对低能尾的效率曲线。
因此即使本报告不再删事件，输入样本的低端形状也可能已经被上游选择雕刻。
\end{warningbox}

\begin{conclusionbox}
原始文件身份、行列数和事件键均可复核，基础变量对全部 4500 个事件有效。
数据足以研究本样本上的相对谱宽改善，但单 run、非均匀空间覆盖和未知 \(E>2\,\mathrm{MeV}\) 预选
使其不足以给出无偏的最终物理分辨率。
\end{conclusionbox}

\section{评价指标：怎样避免“看起来更窄”的假象}
\label{sec:metrics}

\begin{questionbox}
为什么主结果使用 \(\Rsixtyeight\) 和 \(\Rninety\)？它们与 Gaussian 核心宽度 \(\sigma/\mu\) 有什么区别？改善量和不确定度如何定义？
\end{questionbox}

\subsection{分位数与全事件宽度}

对随机变量 \(X\)，第 \(p\) 分位数定义为

\begin{equation}
Q_p(X)
=
\inf\left\{x:\Pr(X\le x)\ge p\right\}.
\label{eq:quantile}
\end{equation}

对归一化相对能量 \(\Erel\)，主指标定义为

\begin{equation}
\Rsixtyeight
=
\frac{Q_{0.84135}(\Erel)-Q_{0.15865}(\Erel)}{2Q_{0.50}(\Erel)},
\label{eq:r68}
\end{equation}

\begin{equation}
\Rninety
=
\frac{Q_{0.95}(\Erel)-Q_{0.05}(\Erel)}{2Q_{0.50}(\Erel)}.
\label{eq:r90}
\end{equation}

两者都使用固定事件集合中的全部有限事件，不需要选择拟合窗口或假定峰形。
\(\Rsixtyeight\) 主要描述中心 68.27\% 的宽度，\(\Rninety\) 对更远的尾部更敏感。

\subsection{Gaussian 核心宽度}

若采用 Gaussian 信号模型，其核心分辨率为

\begin{equation}
R_{\mathrm G}=\frac{\sigma}{\mu}.
\label{eq:gaussian-resolution}
\end{equation}

对理想 Gaussian 分布，

\begin{equation}
\Rsixtyeight\simeq\frac{\sigma}{\mu},
\qquad
\Rninety\simeq1.64485\frac{\sigma}{\mu}.
\label{eq:gaussian-equivalence}
\end{equation}

但真实谱中若有背景、偏斜、长尾或阈值选择，模型可能把真实尾部错误地解释成背景，
从而得到过小的 \(\sigma\)。因此只有峰形模型通过拟合优度检验后，\cref{eq:gaussian-resolution}才允许作为辅助结果。

\begin{warningbox}[title={三个指标不能混用}]
\(\Rsixtyeight\) 和 \(\Rninety\) 是全事件非参数宽度；
\(R_{\mathrm G}=\sigma/\mu\) 是特定峰形模型中的核心宽度。
本报告所有正式比较同时给出 \(\Rsixtyeight\) 和 \(\Rninety\)。
未通过参数 bootstrap GOF 的 \(R_{\mathrm G}\) 只用于说明模型失败，不作为成果。
\end{warningbox}

\subsection{改善量}

对任一宽度 \(R\)，定义绝对变化和相对改善率：

\begin{align}
\Delta R
&=
R_{\mathrm{corr}}-R_{\mathrm{raw}},
\label{eq:delta-r}\\
I_R
&=
1-\frac{R_{\mathrm{corr}}}{R_{\mathrm{raw}}}.
\label{eq:relative-improvement}
\end{align}

\(\Delta R<0\) 表示变窄，单位常写为“百分点”；\(I_R>0\) 表示相对改善。
例如 \(7.893\%\rightarrow6.779\%\) 的绝对变化是 \(-1.115\) 个百分点，
相对改善约 \(14.12\%\)，两种百分数不能混淆。

\begin{conclusionbox}
主结论由固定事件集合上的 \(\Rsixtyeight\) 与 \(\Rninety\) 支撑；
Gaussian 核心宽度必须额外通过峰形 GOF。这样可以避免通过改变窗口、背景函数或事件集合得到虚假的窄峰。
\end{conclusionbox}

\section{全部处理步骤的时间线}
\label{sec:timeline}

\begin{questionbox}
从拿到数据到冻结当前结果，总共做了哪些处理？每一步产生了什么成果，最后是保留、修改还是废弃？
\end{questionbox}

\small
\begin{longtable}{C{0.85cm} p{3.4cm} p{5.3cm} p{2.8cm}}
\caption{完整研究时间线。负结果与被废弃的方法同样保留在记录中。}
\label{tab:timeline}\\
\toprule
阶段 & 做了什么 & 主要观察或成果 & 最终处理 \\
\midrule
\endfirsthead
\toprule
阶段 & 做了什么 & 主要观察或成果 & 最终处理 \\
\midrule
\endhead
1 & 下载并计算 SHA-256 & 文件身份固定；4500 事件、188 字段 & 保留 \\
2 & 审计事件键与全部变量 & 无重复事件；发现 9 个全空列、派生分支无穷值和常数位置列 & 白名单化 \\
3 & 审计 14 个 cut flags & 0/1 语义混合，无法安全解释 & 不加额外 cut \\
4 & 早期派生分支比较 & 得到看似很小的 Gaussian core；发现若干分支几乎仅差常数尺度 & 仅作历史记录 \\
5 & 复核旧峰形评价器 & 结果随窗口大幅变化；正式 GOF 拒绝关键谱形 & 废弃旧核心结论 \\
6 & 冻结独立协议 & 只允许 8 个基础字段，不读已有校正量 & 保留 \\
7 & 第一版基础变量原型 & \(\Rsixtyeight\) 略改善但 \(\Rninety\) 变差；折间方法跳动 & 废弃 \\
8 & 建立分通道物理消融矩阵 & 单独修正 S1 或 S2 都变差；同时修正才稳定改善 & A4 主方案 \\
9 & 增加线性 file 漂移 & 一种拆分略好、另一种变差，无稳定增益 & 不采用 A5 \\
10 & 建立完整基础变量方法矩阵 & \(dt+r^2\) 后做 \(x/y\) 残差面取得更大改善 & challenger \\
11 & 两种文件级嵌套五折 & 交错文件和连续文件块均改善；全部事件获得 OOF 预测 & 保留 \\
12 & 1000 次配对 file bootstrap & A4 和 challenger 的 \(\Delta\Rsixtyeight\) 区间均不跨零 & 条件证据 \\
13 & 多初值、Poisson deviance 与参数 bootstrap GOF & A4 核心模型被拒；challenger 仅连续块通过 & 主报 R68/R90 \\
14 & 冻结方案和下一批数据需求 & A4 定为主方案，challenger 定为候选；要求无阈值数据和独立 run & 待盲测 \\
\bottomrule
\end{longtable}
\normalsize

\begin{conclusionbox}
当前结果不是一次性试出最好数字，而是经过“发现问题—否决旧指标—冻结规则—保留负结果—双拆分验证—模型审计”的完整链条。
后续各章按时间顺序展开\cref{tab:timeline}中的关键步骤。
\end{conclusionbox}

\section{早期派生分支探索：为什么看似更好的结果被作废}

\begin{questionbox}
最初为什么使用已有派生分支？得到了哪些漂亮数字？后来发现了什么问题，为什么这些数字不能继续作为结论？
\end{questionbox}

\subsection{步骤 4A：为什么尝试派生分支}

最初的自然想法是比较文件中已经存在的 corrected、lifetime、ub/stretch 和 MCPAF 分支，
把它们当作 benchmark，快速了解 S1--S2 组合可能达到的宽度。这一阶段发生在独立性规则冻结之前，
因此后来只保留为历史对照，不参与 A4 或 challenger 的设计选择。

历史单次测试拆分为 2672 个训练事件、891 个验证事件和 937 个原始测试事件；
由于派生分支存在无穷值，实际共同测试事件为 936。

\subsection{步骤 4B：观察到的早期数字}

旧 Gaussian 加指数背景拟合曾给出\cref{tab:historical-branches}中的核心宽度。

\begin{table}[H]
\centering
\caption{早期派生分支的单次测试结果。这些数值来自旧峰形模型，仅作为历史记录。}
\label{tab:historical-branches}
\begin{tabularx}{0.96\textwidth}{Y R{3.0cm}}
\toprule
分支组合 & 旧 core \(\sigma/\mu\) \\
\midrule
\texttt{qS1ub\_C + qS2Bdesub\_C\_stretch} & 4.748\% \\
\texttt{qS1ub\_C + qS2Bub\_C\_stretch} & 4.748\% \\
MCPAF desaturated & 5.590\% \\
\texttt{qS1C\_max + qS2Belife\_max} & 6.262\% \\
\texttt{qS1\_max + qS2B\_max} & 6.453\% \\
S2B only & 8.200\% \\
S1 only & 8.638\% \\
\bottomrule
\end{tabularx}
\end{table}

旧派生分支嵌套分析还曾报告 documented corrected/lifetime 的核心宽度
\(6.443\%\rightarrow5.021\%\)，以及 ub/stretch 的
\(4.820\%\rightarrow4.695\%\)。这些数字当时看起来很有吸引力。

\subsection{步骤 4C：分支等价性检查发现的信息不足}

事件级比值检查得到

\begin{equation}
\frac{\texttt{qS2Belife\_max}}{\texttt{qS2BC\_max}}
=1.064133598,
\qquad
\operatorname{std}=2.85\times10^{-8}.
\label{eq:branch-ratio-lifetime}
\end{equation}

这说明两者在当前样本中几乎只差一个常数尺度，不能据此证明发生了事件级电子寿命修正。
另一组 stretch 分支满足

\begin{equation}
\frac{\texttt{qS2Bdesub\_C\_stretch}}
{\texttt{qS2Bub\_C\_stretch}}
\simeq0.999991,
\label{eq:branch-ratio-stretch}
\end{equation}

同样几乎等价。分支名称不能替代数值审计。

\subsection{步骤 5A：窗口系统学暴露模型脆弱性}

documented baseline 的旧核心宽度随窗口变化为：

\begin{center}
\begin{tabular}{cc}
\toprule
拟合窗口 & 旧 core \(\sigma/\mu\) \\
\midrule
\([0.76,1.24]\) & 6.856\% \\
\([0.78,1.22]\) & 6.262\% \\
\([0.80,1.20]\) & 5.563\% \\
\([0.82,1.18]\) & 5.041\% \\
\bottomrule
\end{tabular}
\end{center}

窗口越窄，模型越容易把尾部排除或解释为背景，核心宽度随之显著变小。
这表明 \(\sigma/\mu\) 主要受建模选择影响，而不是一个稳定的物理量。

\subsection{步骤 5B：正式 GOF 否决旧评价器}

对 ub/stretch corrected 谱使用固定窗口 Gaussian 加指数背景、40-bin Poisson deviance，
并进行 200 次参数 bootstrap 后得到

\begin{equation}
p_{\mathrm{boot}}=\frac{1}{201}=0.004975.
\label{eq:historical-gof}
\end{equation}

在 \(0.05\) 显著性水平下明确拒绝该峰形模型。documented corrected 谱在三种背景下的渐近诊断也均给出很小的 \(p\) 值。

\begin{warningbox}[title={为什么旧的窄峰不能保留为成果}]
如果模型不能描述观测谱，\(\sigma\) 就只是错误模型内部的参数。
旧的 \(6.443\%\rightarrow5.021\%\) 和 \(4.820\%\rightarrow4.695\%\)
只能记录为“曾经观察到但已被审计否决的历史结果”，不能称为能量分辨率改善。
\end{warningbox}

\begin{conclusionbox}
本阶段最大的成果不是得到一个更小数字，而是发现旧分支信息不足、旧峰形模型会产生虚假窄核心，
并据此把主评价指标改为全事件 \(\Rsixtyeight\) 与 \(\Rninety\)。
\end{conclusionbox}

\section{第一版基础变量原型：一次必要的失败}

\begin{questionbox}
冻结变量白名单后，第一版独立方法做了什么？为什么中心宽度略有改善却仍被判定失败？
\end{questionbox}

\subsection{步骤 6：冻结独立协议}

在历史评价器被否决后，首先写下 \path{INDEPENDENT_PROTOCOL.md}，固定\cref{eq:whitelist}中的输入字段，
要求响应图和权重只由训练侧学习，测试文件不得参与调参。这一步把后续探索与已有修正分支完全隔开。

\subsection{步骤 7A：原型方法}

第一版 \path{independent_reconstruction.py} 分别为 S1 和 S2B 尝试
none、\(dt\)、\(r^2\)、\(dt+r^2\) 等粗分箱响应修正，再在内层选择组合权重。
其思想方向正确，但仍借用了旧 Gaussian 加指数拟合来选择候选，而且每个外折只有一次简化内层选择。

不同外折选择的通道修正从“两个都不修正”跳到“S1 只修正 \(dt\)”或“S2 修正 \(dt+r^2\)”，
正权重在约 0.23--0.77 之间明显跳动，说明 4500 个非均匀事件不足以稳定支持当时的候选空间。

\subsection{步骤 7B：多指标结果与否决理由}

原型在全事件评价下得到

\begin{align}
\Rsixtyeight &: 8.908\%\longrightarrow8.519\%,
\label{eq:prototype-r68}\\
\Rninety &: 16.176\%\longrightarrow16.215\%.
\label{eq:prototype-r90}
\end{align}

中心 68\% 略变窄，但 90\% 宽度略变差；旧 core 指标也从 \(6.585\%\) 变为 \(6.789\%\)。
300 次冻结 OOF bootstrap 的旧 core 差值区间为

\begin{equation}
[-0.143,+1.476]\ \text{个百分点},
\label{eq:prototype-bootstrap}
\end{equation}

跨过零。结论随指标改变，且折间方案不稳定，因此该原型不满足“多指标、双拆分、稳定选择”的接受标准。

\begin{conclusionbox}
第一版原型证明了“只用基础变量”在技术上可行，也暴露了三个问题：不能再用旧峰形选择模型；
必须减少自由度；必须同时审查中心和尾部。它被保留为失败路径，而不是被悄悄删除。
\end{conclusionbox}

\section{响应校正与 OOF 重建的一般方法}

\begin{questionbox}
响应图在数学上做什么？如何防止模型直接把每个事件“拉回常数”？训练、内层验证与外层测试分别承担什么角色？
\end{questionbox}

\subsection{局部相对响应}

对观测通道 \(A\) 和响应坐标 \(u\)，只在外折训练数据中把 \(u\) 分为少量等频箱，
在每个箱中估计鲁棒信号中心 \(C_A(u)\)。相对于全局中心的响应为

\begin{equation}
\rho_A(u)
=
\frac{C_A(u)}{C_A^{\mathrm{global}}}.
\label{eq:relative-response}
\end{equation}

把响应模型平滑或低阶参数化后，校正写作

\begin{equation}
A^{\mathrm{corr}}
=
\frac{A^{\mathrm{obs}}}{\rho_A(u)}.
\label{eq:response-correction}
\end{equation}

若依次校正多个坐标，则

\begin{equation}
A^{(k)}
=
\frac{A^{(k-1)}}{\rho_k(u_k)},
\qquad
A^{(0)}=A^{\mathrm{obs}}.
\label{eq:sequential-correction}
\end{equation}

模型拟合的是箱级中位中心，不是每个事件与常数 1 的残差；这降低了追踪单事件统计涨落的风险。

\subsection{训练侧电子寿命模型}

S2B 的训练侧中心随漂移时间拟合为

\begin{align}
\ln C_{\StwoB}(t_{\mathrm d})
&=a+b(t_{\mathrm d}-t_0),
\qquad b\le0,
\label{eq:lifetime-fit}\\
\StwoB^{\mathrm{corr}}
&=\StwoB^{\mathrm{obs}}
\exp\left[-b(t_{\mathrm d}-t_0)\right].
\label{eq:lifetime-correction}
\end{align}

当 \(b<0\) 时，表观寿命为

\begin{equation}
\tau_e=-\frac{1}{b}.
\label{eq:tau-from-slope}
\end{equation}

当前外折得到约 \(1.214\)--\(1.349\,\mathrm{ms}\)，但由于源位置与漂移时间相关，
它只能视为校正参数，不是正式电子寿命测量。

\subsection{嵌套选择与 OOF 预测}

对第 \(k\) 个外折，方法、响应参数和组合权重都只由其余四折学习。
内层验证选择超参数 \(\theta\)：

\begin{equation}
\widehat{\theta}_k
=
\argminop_{\theta\in\Theta}
\Rsixtyeight\!\left(
\mathcal D_{k}^{\mathrm{inner\text{-}val}};
\theta
\right).
\label{eq:inner-selection}
\end{equation}

事件 \(i\) 的 OOF 预测定义为

\begin{equation}
\Erel_i^{\OOF}
=
\mathcal M^{(-k(i))}(\symbf{x}_i),
\label{eq:oof}
\end{equation}

其中 \(\mathcal M^{(-k)}\) 从未看过第 \(k\) 外折。所有最终宽度由 \(4500\) 个 OOF 预测合并后计算。

\begin{conclusionbox}
低自由度箱级响应、训练侧拟合、内层调参和外层 OOF 共同构成防泄漏链条。
测试事件只在最后评价，不参与中心、响应面、方法族或权重选择。
\end{conclusionbox}

\section{物理可解释主方案 A4}

\begin{questionbox}
A4 具体怎样校正两个通道？为什么单独校正 S1 或 S2 会变差，而联合校正反而稳定改善？为什么不加入 file 漂移？
\end{questionbox}

\subsection{步骤 8A：通道消融矩阵}

为了区分每个物理步骤的贡献，建立六个固定候选：

\begin{description}[leftmargin=1.5cm,labelwidth=1.1cm]
  \item[A0] 原始 S1 与原始 S2B；
  \item[A1] 只对 S2B 做指数寿命修正；
  \item[A2] 只对 S1 做 \(dt\) 和 \(r^2\) 响应修正；
  \item[A3] 只对 S2B 做寿命与 \(r^2\) 修正；
  \item[A4] 同时修正 S1 的 \((dt,r^2)\) 和 S2B 的 \((\text{lifetime},r^2)\)；
  \item[A5] 在 A4 后再对两个通道加入线性 file 漂移。
\end{description}

每个局部响应只使用五个等频箱。修正后的通道按训练中心归一化：

\begin{align}
s_1
&=
\frac{\Sone^{\mathrm{corr}}}{C_{\Sone,\mathrm{train}}},
\label{eq:s1-normalized}\\
s_2
&=
\frac{\StwoB^{\mathrm{corr}}}{C_{\StwoB,\mathrm{train}}}.
\label{eq:s2-normalized}
\end{align}

组合相对能量为

\begin{equation}
e_{\alpha}
=
\alpha s_1+(1-\alpha)s_2,
\qquad 0.10\le\alpha\le0.90,
\label{eq:a4-energy}
\end{equation}

再用训练侧组合中心归一化：

\begin{equation}
\Erel
=
\frac{e_{\alpha}}{C_{e,\mathrm{train}}}.
\label{eq:relative-energy}
\end{equation}

\(\alpha\) 只在内层验证集选择，外层测试集从不参与。

\subsection{步骤 8B：消融结果}

\begin{table}[H]
\centering
\caption{A0--A5 的全事件 OOF \(\Rsixtyeight\)。单通道修正均恶化，只有联合修正稳定改善。}
\label{tab:a4-ablation}
\begin{tabular}{lcc}
\toprule
候选 & file mod 5 & 连续 file 块 \\
\midrule
A0 raw & 7.774\% & 7.893\% \\
A1 S2B lifetime & 9.351\% & 9.383\% \\
A2 S1 \((dt,r^2)\) & 9.146\% & 9.085\% \\
A3 S2B \((\text{lifetime},r^2)\) & 8.532\% & 8.522\% \\
\rowcolor{SoftTeal}
A4 both & \textbf{6.773\%} & \textbf{6.779\%} \\
A5 A4 + linear file drift & 6.824\% & 6.773\% \\
\bottomrule
\end{tabular}
\end{table}

单独把一个通道“校平”会打破原始 S1 与 S2B 中部分相互补偿的位置依赖，
因此 A1--A3 比 raw 更差。联合修正后，两个通道被放到一致的响应基准，再由正权重组合，才得到稳定改善。
这一负结果也是为什么最终方案不能简单写成“只做电子寿命修正”。

\subsection{步骤 9：为什么不采用 A5}

线性 file 漂移在 file mod 5 拆分中使 \(\Rsixtyeight\) 从 A4 的 \(6.773\%\) 变为 \(6.824\%\)，
即变差；在连续块中从 \(6.779\%\) 变为 \(6.773\%\)，只改善约 \(0.006\) 个百分点。
增益不跨拆分稳定，且 fileNumber 本身只是采集顺序代理，不是物理坐标，因此 A5 不被采用。

\subsection{A4 的响应平坦度}

在 file mod 5 OOF 诊断中，漂移时间分箱峰中心的峰峰变化从 \(4.457\%\) 降至 \(3.047\%\)，
径向 \(r^2\) 分箱峰中心的峰峰变化从 \(10.052\%\) 降至 \(5.022\%\)。

\begin{figure}[H]
\centering
\includegraphics[width=0.98\textwidth]{channel_flatness.png}
\caption{A4 修正前后的 OOF 响应平坦度。联合通道修正显著减小径向依赖，并降低漂移时间依赖；
图中点只用于事后诊断，不参与 OOF 模型训练。}
\label{fig:channel-flatness}
\end{figure}

\begin{conclusionbox}
A4 是当前主方案，因为它以低自由度、物理可解释的方式同时修正两个通道，
在两种文件拆分中都改善 \(\Rsixtyeight\)、\(\Rninety\) 和 \(\pm20\%\) 外事件比例。
单通道方案和额外 file 漂移作为明确负结果保留。
\end{conclusionbox}

\section{更灵活的组合响应 challenger}

\begin{questionbox}
challenger 与 A4 有什么根本区别？它为什么能得到更大的改善？为什么暂时只能叫候选，而不能直接替代主方案？
\end{questionbox}

\subsection{步骤 10A：原始组合与正权重}

challenger 先在训练折中构造

\begin{equation}
e_{\alpha}^{(0)}
=
\frac{\StwoB}{\median_{\mathrm{train}}(\StwoB)}
+
\alpha
\frac{\Sone}{\median_{\mathrm{train}}(\Sone)},
\qquad \alpha>0.
\label{eq:challenger-raw}
\end{equation}

代码在每个外折内部比较 17 个正权重候选，避免不物理的负通道系数。
由于\cref{eq:challenger-raw}与 A4 的权重参数化不同，两条管线的 raw 宽度不同；
只能比较同一管线同一行的修正前后，不能把 A4 raw 与 challenger raw 直接相减。

\subsection{步骤 10B：箱级低阶响应面}

对任意坐标 \(x_j\)，先用训练侧位置和尺度标准化：

\begin{equation}
u_j
=
\frac{x_j-\median_{\mathrm{train}}(x_j)}{s_j}.
\label{eq:standard-coordinate}
\end{equation}

联合 \(dt+r^2\) 二阶响应面为

\begin{equation}
f_{t,r^2}(u,v)
=
\beta_0+\beta_1u+\beta_2v
+\beta_3u^2+\beta_4v^2+\beta_5uv.
\label{eq:dt-r2-surface}
\end{equation}

随后对残余 \(x/y\) 依赖拟合

\begin{equation}
f_{x,y}(u_x,u_y)
=
\gamma_0+\gamma_1u_x+\gamma_2u_y
+\gamma_3u_x^2+\gamma_4u_xu_y+\gamma_5u_y^2.
\label{eq:xy-surface}
\end{equation}

模型先把训练事件分为粗空间箱，只拟合各箱 \(\ln e\) 的中位数 \(y_c\)。
带箱内事件数权重和岭正则的目标可写为

\begin{equation}
\widehat{\symbf{\beta}}
=
\argminop_{\symbf{\beta}}
\left[
\sum_c n_c\left(y_c-\symbf{x}_c^{\mathsf T}\symbf{\beta}\right)^2
+\lambda\lVert\symbf{\beta}\rVert_2^2
\right].
\label{eq:cell-fit}
\end{equation}

校正应用为

\begin{equation}
e^{\mathrm{corr}}
=
e^{(0)}
\exp\left[
-\operatorname{clip}
\left(f(\symbf{u})-f_0,-0.30,0.30\right)
\right].
\label{eq:surface-correction}
\end{equation}

\(\pm0.30\) 的对数裁剪限制单次响应修正，避免稀疏区域出现无界外推。

\subsection{步骤 10C：完整候选矩阵}

\begin{table}[H]
\centering
\caption{challenger 方法矩阵的全事件 OOF \(\Rsixtyeight\)。负结果没有被过滤。}
\label{tab:challenger-matrix}
\begin{tabular}{lcc}
\toprule
响应方法 & file mod 5 & 连续 file 块 \\
\midrule
raw & 8.019\% & 8.125\% \\
\(dt\) quadratic & 7.553\% & 7.584\% \\
\(r^2\) quadratic & 6.389\% & 6.339\% \\
joint \(dt+r^2\) & 6.283\% & 6.221\% \\
\(dt\rightarrow r^2\) & 6.383\% & 6.318\% \\
\(r^2\rightarrow dt\) & 6.138\% & 6.172\% \\
\(x/y\) quadratic & 6.734\% & 6.685\% \\
\(dt\rightarrow x/y\) & 6.640\% & 6.615\% \\
\rowcolor{SoftGold}
joint \(dt+r^2\rightarrow x/y\) & \textbf{5.833\%} & \textbf{5.973\%} \\
\bottomrule
\end{tabular}
\end{table}

内层 nested winner 在两种拆分的全部十个外折中都选择
\texttt{dt\_r2\_then\_xy}。这说明当前数据上联合深度--半径响应后仍存在横向残差。

\subsection{为什么 challenger 仍不是最终方案}

challenger 直接校正组合能量，包含两个二阶响应面，比 A4 自由度更高。
外源样本的空间分布、未知 \(E>2\,\mathrm{MeV}\) 预选和真实能谱成分可能彼此相关，
更灵活的模型可能把源分布或选择效应的一部分吸收到“响应校正”中。
内部 OOF 可以防止最直接的事件泄漏，但不能替代第二个 run 的冻结盲测。

\begin{conclusionbox}
challenger 是当前数据上表现最强的候选，在全部十个外折中胜过自己的 raw 基线；
但它的解释性和迁移风险弱于 A4。因此 A4 作为物理主方案，challenger 作为冻结候选，
二者都必须进入下一批独立数据测试。
\end{conclusionbox}

\section{文件级嵌套交叉验证设计}

\begin{questionbox}
为什么不能随机拆分事件？两种外层拆分分别检验什么？怎样保证测试文件没有参与方法与权重选择？
\end{questionbox}

\subsection{步骤 11A：为什么按 file 分组}

同一 file 中的事件共享近似的探测器状态、采集时间和处理条件。
若随机拆分事件，同一 file 的模式会同时出现在训练和测试中，使测试过于乐观。
因此本研究的最小独立单元是 file，而不是 event。

\subsection{步骤 11B：两套外层五折}

第一套为交错文件拆分：

\begin{equation}
k=\texttt{fileNumber}\bmod5.
\label{eq:file-mod-fold}
\end{equation}

它主要检验对不同文件之间的插值能力。第二套把排序后的 985 个唯一 fileNumber 分成五个连续块，
块内事件数依次为 869、939、928、873、891；它更接近对相邻时间区间的外推测试，
因此被用作主报告拆分。

\subsection{步骤 11C：内层选择规则}

外折训练部分再次按 file 分组做内部验证。A4 选择 \(\alpha\)；challenger 同时选择正权重和响应方法族。
候选方法若在最佳 \(\Rsixtyeight\) 的 \(0.05\) 个百分点以内，则优先选择自由度最低者。
外层测试只被预测一次。

\subsection{五个连续测试块的稳定性}

\begin{table}[H]
\centering
\caption{连续文件块逐折 \(\Rsixtyeight\)。A4 和 challenger 在五个未见时间块中全部优于各自 raw。}
\label{tab:fold-results}
\begin{tabular}{c cc cc}
\toprule
外折 & A4 raw & A4 corrected & challenger raw & challenger corrected \\
\midrule
0 & 8.168\% & 6.995\% & 8.412\% & 5.726\% \\
1 & 7.593\% & 6.801\% & 7.891\% & 6.079\% \\
2 & 7.904\% & 6.736\% & 7.895\% & 6.287\% \\
3 & 7.979\% & 6.658\% & 8.037\% & 5.634\% \\
4 & 7.948\% & 6.685\% & 8.278\% & 5.989\% \\
\bottomrule
\end{tabular}
\end{table}

\begin{conclusionbox}
两套按文件五折回答不同的泛化问题。最严格的连续块结果不是由某一个外折驱动：
A4 和 challenger 的五个测试块全部改善。这比只报告合并后的一个数字更有说服力。
\end{conclusionbox}

\section{主要 OOF 结果}
\label{sec:main-results}

\begin{questionbox}
最终主方案和候选方案分别把 \(\Rsixtyeight\)、\(\Rninety\) 改善到多少？改善是绝对百分点还是相对比例？事件保留情况如何？
\end{questionbox}

\begin{table}[H]
\centering
\caption{最终全事件 OOF 结果。每一行只比较同一管线自己的 raw 与 corrected。}
\label{tab:main-results}
\begin{tabular}{llcccc}
\toprule
方案 & 外层拆分 & Raw \(\Rsixtyeight\) & Corr. \(\Rsixtyeight\) & Raw \(\Rninety\) & Corr. \(\Rninety\) \\
\midrule
A4 & file mod 5 & 7.774\% & 6.773\% & 14.740\% & 13.262\% \\
\rowcolor{SoftTeal}
A4 & 连续 file 块 & 7.893\% & \textbf{6.779\%} & 14.783\% & \textbf{13.278\%} \\
challenger & file mod 5 & 8.019\% & 5.833\% & 15.127\% & 11.776\% \\
\rowcolor{SoftGold}
challenger & 连续 file 块 & 8.125\% & \textbf{5.973\%} & 15.029\% & \textbf{11.826\%} \\
\bottomrule
\end{tabular}
\end{table}

\subsection{A4 主方案的量化改善}

连续块中，

\begin{align}
\Delta\Rsixtyeight
&=6.779\%-7.893\%=-1.115\ \text{个百分点},
\label{eq:a4-delta-r68}\\
I_{R_{68}}
&=1-\frac{6.779}{7.893}=14.12\%,
\label{eq:a4-relative-r68}\\
\Delta\Rninety
&=13.278\%-14.783\%=-1.505\ \text{个百分点}.
\label{eq:a4-delta-r90}
\end{align}

\(\pm20\%\) 窗口外事件比例同时从 \(3.756\%\) 降至 \(2.333\%\)。
file mod 5 拆分也得到一致方向：\(\Rsixtyeight\) 绝对降低 \(1.001\) 个百分点，
\(\Rninety\) 绝对降低 \(1.477\) 个百分点。

\subsection{challenger 的量化改善}

连续块中，

\begin{align}
\Delta\Rsixtyeight
&=5.973\%-8.125\%=-2.152\ \text{个百分点},
\label{eq:challenger-delta-r68}\\
I_{R_{68}}
&=1-\frac{5.973}{8.125}=26.48\%,
\label{eq:challenger-relative-r68}\\
\Delta\Rninety
&=11.826\%-15.029\%=-3.202\ \text{个百分点}.
\label{eq:challenger-delta-r90}
\end{align}

file mod 5 中的 \(\Rsixtyeight\) 相对改善约为 \(27.26\%\)。

\begin{figure}[H]
\centering
\includegraphics[width=0.98\textwidth]{result_summary.png}
\caption{四组全事件 OOF 宽度比较。A4 与 challenger 在交错文件和连续文件块两种拆分中均同时改善 \(\Rsixtyeight\) 和 \(\Rninety\)。}
\label{fig:result-summary}
\end{figure}

\begin{figure}[H]
\centering
\includegraphics[width=0.98\textwidth]{oof_spectra.png}
\caption{A4 与 challenger 的 OOF 相对能量分布。图只用于展示全体分布形状；正式宽度来自分位数，
没有通过选择拟合窗口或删除尾部事件来计算。}
\label{fig:oof-spectra}
\end{figure}

所有四组比较都保留 \(4500/4500\) 个事件。需要注意，A4 中相对于变窄后 \(2\Rsixtyeight\) 的尾部比例略有上升，
因此只说“中心变窄”不够；本报告同时展示 \(\Rninety\) 和固定 \(\pm20\%\) 尾部，后两者均改善。

\begin{resultbox}
主成果是 A4 连续块 \(\Rsixtyeight:7.893\%\rightarrow6.779\%\)，
候选成果是 challenger 连续块 \(8.125\%\rightarrow5.973\%\)。
两者都来自 4500 个 OOF 事件且没有额外事件 cut；A4 是主方案，challenger 是待独立 run 验证的候选。
\end{resultbox}

\section{Bootstrap 不确定度与其适用范围}

\begin{questionbox}
改善是否可能只是有限样本波动？为什么按 file 重采样而不是按 event？当前区间为什么仍不是完整流程不确定度？
\end{questionbox}

\subsection{步骤 12：配对 file bootstrap}

设唯一文件集合为 \(\{f_1,\ldots,f_F\}\)，每次有放回抽取同样数量的 file：

\begin{equation}
f_1^{\ast},\ldots,f_F^{\ast}
\overset{\mathrm{i.i.d.}}{\sim}
\operatorname{Unif}\{f_1,\ldots,f_F\}.
\label{eq:file-bootstrap}
\end{equation}

同一事件的 raw 和 corrected OOF 预测始终成对出现。第 \(b\) 次重采样计算

\begin{equation}
\Delta\Rsixtyeight^{\ast b}
=
R_{68,\mathrm{corr}}^{\ast b}
-
R_{68,\mathrm{raw}}^{\ast b}.
\label{eq:bootstrap-delta}
\end{equation}

百分位区间为

\begin{equation}
\mathrm{CI}_{95\%}
=
\left[
Q_{0.025}\!\left(\Delta R^{\ast}\right),
Q_{0.975}\!\left(\Delta R^{\ast}\right)
\right].
\label{eq:bootstrap-ci}
\end{equation}

\subsection{1000 次重采样结果}

\begin{table}[H]
\centering
\caption{冻结 OOF 预测上的配对 file bootstrap。区间单位为百分点。}
\label{tab:bootstrap-results}
\begin{tabular}{llcc}
\toprule
方案 & 外层拆分 & \(\Delta\Rsixtyeight\) 中位数 & 95\% 区间 \\
\midrule
A4 & file mod 5 & $-1.027$ & $[-1.257,-0.793]$ \\
A4 & 连续 file 块 & $-1.125$ & $[-1.345,-0.898]$ \\
challenger & file mod 5 & $-2.206$ & $[-2.506,-1.899]$ \\
challenger & 连续 file 块 & $-2.192$ & $[-2.488,-1.898]$ \\
\bottomrule
\end{tabular}
\end{table}

四个区间均不跨零，说明在当前冻结预测和文件重采样定义下，改善不太可能由少数 file 的随机波动单独造成。

\begin{warningbox}[title={这是条件区间，不是完整流程误差}]
每次 bootstrap 只重采样已经生成的 OOF 预测，没有重新学习响应图、方法族和权重。
因此\cref{tab:bootstrap-results}条件于当前冻结模型，未包含模型选择不确定度。
相邻 file 也可能相关。下一阶段应在完整管线上做连续块 bootstrap 或在新 run 上盲测。
\end{warningbox}

\begin{conclusionbox}
按 file 配对重采样比事件级 bootstrap 更尊重数据的相关结构，且当前改善区间均不跨零；
但它回答的是“冻结预测是否稳定”，不是“从头重跑整个分析会有多大误差”。
\end{conclusionbox}

\section{峰形拟合与拟合优度审计}

\begin{questionbox}
如果仍想估计 Gaussian 核心宽度，怎样检验峰形模型？为什么 A4 的 3.7\%--4.0\% 不能作为结果？challenger 为什么也只通过了一种拆分？
\end{questionbox}

\subsection{步骤 13A：固定窗口混合模型}

审计使用固定相对窗口 \(0.8\le x\le1.2\)，模型为截断 Gaussian 信号与正定背景的混合：

\begin{equation}
p(x)
=
f_{\mathrm s}p_{\mathrm G}^{\mathrm{tr}}(x\mid\mu,\sigma)
+(1-f_{\mathrm s})p_{\mathrm B}(x\mid\symbf{\theta}).
\label{eq:mixture-model}
\end{equation}

截断 Gaussian 写为

\begin{equation}
p_{\mathrm G}^{\mathrm{tr}}(x\mid\mu,\sigma)
=
\frac{
\phi\!\left((x-\mu)/\sigma\right)/\sigma
}{
\Phi\!\left((b-\mu)/\sigma\right)
-\Phi\!\left((a-\mu)/\sigma\right)
},
\qquad a=0.8,\ b=1.2,
\label{eq:truncated-gaussian}
\end{equation}

背景依次测试指数、一阶 Bernstein 和二阶 Bernstein。每次正式拟合使用 12 个初值，
检查优化成功、多初值一致和参数是否贴边。

\subsection{步骤 13B：Poisson deviance 与参数 bootstrap}

把固定窗口分为 40 个箱，观测数为 \(n_b\)，模型期望为 \(\lambda_b\)。Poisson deviance 为

\begin{equation}
D
=
2\sum_b
\left[
\lambda_b-n_b+n_b\ln\left(\frac{n_b}{\lambda_b}\right)
\right],
\label{eq:poisson-deviance}
\end{equation}

其中 \(n_b=0\) 时对数项定义为零。随后从拟合模型模拟 \(B=200\) 个 toy 谱，
每个 toy 都重新拟合。参数 bootstrap 的 \(p\) 值为

\begin{equation}
p_{\mathrm{boot}}
=
\frac{1+\sum_{j=1}^{B}\mathbb I(D_j^{\ast}\ge D_{\mathrm{obs}})}{B+1}.
\label{eq:gof-pvalue}
\end{equation}

接受标准为 \(p_{\mathrm{boot}}>0.05\)，同时无边界退化且多初值稳定。

\subsection{步骤 13C：最终 GOF 结果}

\begin{table}[H]
\centering
\caption{Gaussian 加二阶 Bernstein 背景的参数 bootstrap GOF。被拒绝的 core 只用于说明模型失败。}
\label{tab:gof-results}
\begin{tabular}{llccc}
\toprule
方案 & 拆分 & 表面 core \(\sigma/\mu\) & \(p_{\mathrm{boot}}\) & 判定 \\
\midrule
A4 & file mod 5 & 4.008\% & 0.004975 & 拒绝 \\
A4 & 连续 file 块 & 3.697\% & 0.004975 & 拒绝 \\
challenger & file mod 5 & 4.389\% & 0.02985 & 拒绝 \\
challenger & 连续 file 块 & 4.045\% & 0.1692 & 通过 \\
\bottomrule
\end{tabular}
\end{table}

虽然 challenger 连续块的单个谱通过 GOF，但相同算法的 file mod 5 谱没有通过；
如果只挑选通过且数字较小的一行，会形成事后选择偏差。因此统一主结论仍使用全事件 \(\Rsixtyeight\) 与 \(\Rninety\)。

\begin{warningbox}
A4 的 \(3.697\%\) 和 \(4.008\%\) 看起来比全事件 \(\Rsixtyeight\) 更好，
但对应峰形模型均被 \(p=0.004975\) 明确拒绝。这两个数字不得写入摘要、结论或对外汇报中的“最终分辨率”。
\end{warningbox}

\begin{conclusionbox}
峰形审计成功阻止了第二次“虚假窄核心”。当前可复核、跨模型假设稳定的结论是
\cref{tab:main-results}中的全事件分位数宽度，而不是某个被拒绝模型的 \(\sigma/\mu\)。
\end{conclusionbox}

\section{负结果总表与从中学到的内容}

\begin{questionbox}
哪些方法没有成功？失败分别说明了什么？这些信息怎样影响最终方案？
\end{questionbox}

\begin{longtable}{p{3.2cm} p{4.0cm} p{4.9cm}}
\caption{保留的负结果。每项均按“观察—否决原因—学习”记录。}
\label{tab:negative-results}\\
\toprule
尝试 & 观察与否决原因 & 学到的内容 \\
\midrule
\endfirsthead
\toprule
尝试 & 观察与否决原因 & 学到的内容 \\
\midrule
\endhead
旧 Gaussian+指数峰形 & core 随窗口从 6.856\% 变到 5.041\%；关键谱 GOF \(p=0.004975\) & 峰形假设必须通过 GOF，主指标改为全事件分位数 \\
派生 corrected/lifetime 名称 & 两个 lifetime 分支事件级比值几乎为常数 & 变量名不能证明修正内容，必须做数值等价性检查 \\
第一版基础变量原型 & \(\Rsixtyeight\) 略好但 \(\Rninety\) 变差，折间方案跳动 & 候选要低自由度，接受标准必须多指标 \\
只做 S2B 寿命（A1） & 连续块 \(\Rsixtyeight=9.383\%\)，比 raw 更差 & 原始两个通道存在补偿，不能单边校平 \\
只做 S1（A2） & 连续块 \(9.085\%\)，比 raw 更差 & 需要同步处理 S1 与 S2B \\
只做 S2B 寿命和半径（A3） & 连续块 \(8.522\%\)，仍比 raw 更差 & 增加径向修正仍不能替代双通道一致校正 \\
线性 file 漂移（A5） & 一种拆分变差，另一种仅好 0.006 个百分点 & 不使用不稳定的时间代理修正 \\
只做 \(dt\) challenger & 连续块从 8.125\% 到 7.584\%，收益有限 & 当前样本的主要剩余结构与横向位置相关 \\
参数化 A4 core & 3.697\%--4.008\% 很小，但两种拆分 GOF 均拒绝 & 不把“拟合成功”误当“模型正确” \\
\bottomrule
\end{longtable}

\begin{conclusionbox}
负结果不是项目的附属品，而是最终方法可信度的一部分：它们排除了单边校正、file 漂移和错误峰形，
并促成了双通道 A4、完整 challenger 矩阵和严格 GOF 审计。
\end{conclusionbox}

\section{最终完成了什么、还没有完成什么}
\label{sec:final}

\begin{questionbox}
到目前为止，哪些工作已经真正完成？哪些结论仍需要新数据？应该如何准确地向导师或合作者汇报？
\end{questionbox}

\subsection{已经完成的工作}

\begin{enumerate}[label=\arabic*.]
  \item 完成原始文件下载、文件哈希、行列数、事件身份和 188 个字段的可复核审计。
  \item 冻结了只使用八个基础字段的独立分析协议，不依赖以前的重建代码或已校正信号分支。
  \item 在基础变量全部有效的条件下保留 \(4500/4500\) 个事件，没有通过额外 cut 人为缩窄谱。
  \item 建立按 file 分组的嵌套 OOF 验证，避免同一 file 同时进入训练与测试。
  \item 建立物理可解释的 A4 双通道修正，在五个连续测试块中全部改善：
  \[
  \Rsixtyeight:7.893\%\longrightarrow6.779\%.
  \]
  \item 建立更强但更灵活的 challenger，在五个连续测试块中同样全部改善：
  \[
  \Rsixtyeight:8.125\%\longrightarrow5.973\%.
  \]
  \item 同时报告 \(\Rninety\)、固定尾部和逐折稳定性，没有只展示最有利指标。
  \item 建立 1000 次配对 file bootstrap；四组 \(\Delta\Rsixtyeight\) 条件区间均不跨零。
  \item 建立固定窗口、三种背景、多初值、参数边界、Poisson deviance 和参数 bootstrap 的峰形审计器。
  \item 发现并废弃旧峰形评价器、单通道修正、第一版不稳定原型和没有稳定收益的 file 漂移。
\end{enumerate}

\subsection{目前尚未证明的内容}

\begin{enumerate}[label=\arabic*.]
  \item 没有完成以 MeV 为单位的绝对能量标定，也没有独立测定 \(g_1\) 和 \(g_{2,\mathrm B}\)。
  \item 没有确定 run 10972 的源、目标峰能量和完整上游选择定义。
  \item 文件已受未知 \texttt{Egt2MeV} 预选，无法完整评估低能尾、效率和选择偏差。
  \item 只有一个 run，尚未验证跨 run、跨源位置、跨能量和长期稳定性。
  \item 当前 bootstrap 没有在每次重采样中重跑完整训练与模型选择。
  \item challenger 尚未在完全独立数据上冻结盲测，因此不能称为最终物理重建。
\end{enumerate}

\begin{resultbox}[title={可以直接对导师或合作者使用的结论}]
在 run 10972 的未知 \(E>2\,\mathrm{MeV}\) 预选样本上，
仅使用原始 S1、底部 S2B、漂移时间和原始 S2 位置建立的独立重建，
在两种按文件划分的 OOF 验证中均降低了全事件非参数宽度。
物理可解释主方案在连续块中将 \(\Rsixtyeight\) 从 \(7.893\%\) 降至 \(6.779\%\)，
更灵活的 challenger 从 \(8.125\%\) 降至 \(5.973\%\)。
这些结果证明独立基础变量重建存在稳定改进空间，但正式能量分辨率仍需无阈值样本和独立 run 盲测。
\end{resultbox}

\begin{conclusionbox}
本阶段已经成功完成“独立相对能量重建改善”的方法开发与内部泛化验证，
但没有把它包装成超出证据范围的最终物理分辨率。A4 和 challenger 的角色、数字和限制均已冻结。
\end{conclusionbox}

\section{下一阶段的具体工作计划}

\begin{questionbox}
接下来最有价值的工作是什么？需要向数据提供者索取哪些信息和数据？怎样进行真正的冻结盲测？
\end{questionbox}

\subsection{优先级 1：获得无能量阈值预选数据}

最重要的数据请求是同一 run 或相似标定 run 的、未经任何重建能量阈值预选的候选样本，
同时提供上游 selection 的明确逻辑和 cut 方向。这样才能重新评价完整低能尾和效率。

\subsection{优先级 2：独立 run 冻结盲测}

在新 run 到来前冻结以下内容：

\begin{itemize}
  \item A4 的五等频箱、寿命斜率约束、径向响应顺序和 \(\alpha\) 网格；
  \item challenger 的候选方法、17 个正权重、二阶项、岭正则和 \(\pm0.30\) 修正裁剪；
  \item file 分组规则、\(\Rsixtyeight\)、\(\Rninety\)、尾部指标与 GOF 接受标准；
  \item 任何新增方法必须先在开发数据上定义，再一次性进入盲测，不能看完新 run 再改规则。
\end{itemize}

\subsection{优先级 3：补足完整不确定度}

在计算允许时，实施“重采样 file 块—重新学习响应面—重新选择方法和权重—重新生成 OOF”的完整 pipeline bootstrap。
并比较普通 file bootstrap、连续块 bootstrap 和跨 run 差异。

\subsection{优先级 4：绝对能标与物理解释}

确认标定源、峰能、峰组成和 run 条件后，利用多个已知峰联合确定相对线性与绝对能标，
再讨论 \(g_1\)、\(g_2\)、电子寿命和以 MeV 表示的能量分辨率。

\noindent\textbf{\color{PandaTeal}本阶段的行动结论：}
下一步最重要的不是继续在同一 4500 个事件上增加模型复杂度，
而是拿到无阈值样本和独立 run，对已经冻结的 A4 与 challenger 做真正盲测。

\clearpage
\appendix

\section{实际使用变量、单位与用途}

\begin{longtable}{p{3.4cm} p{2.0cm} p{6.2cm}}
\caption{正式独立流程的字段白名单。}
\label{tab:variables}\\
\toprule
字段 & 单位 & 用途 \\
\midrule
\endfirsthead
\toprule
字段 & 单位 & 用途 \\
\midrule
\endhead
\texttt{runNumber} & -- & 数据身份检查；当前固定为 10972 \\
\texttt{fileNumber} & -- & 事件分组、外层拆分和事后时间诊断；challenger 不把它作为响应特征 \\
\texttt{eventNumber} & -- & 与 run、file 共同组成唯一事件键 \\
\texttt{qS1\_max} & p.e. & 主 S1 面积，作为光通道基础观测量 \\
\texttt{qS2B\_max} & p.e. & 主 S2 的底部阵列面积，作为电荷通道基础观测量 \\
\texttt{dt} & ns & S1--S2 漂移时间；分析中除以 1000 转为 \(\mu\mathrm s\) \\
\texttt{xS2T\_max} & mm & 基础 S2 横向位置 x，用于 \(r^2\) 和 challenger 横向残差 \\
\texttt{yS2T\_max} & mm & 基础 S2 横向位置 y，用于 \(r^2\) 和 challenger 横向残差 \\
\bottomrule
\end{longtable}

变量命名参考项目提供的 TechNote\cite{technote}：\texttt{q} 表示面积/电荷，\texttt{t} 表示时间，
\texttt{x/y} 表示重建位置，\texttt{T/B} 表示顶部/底部阵列，\texttt{C/cor} 通常表示已校正。
本报告只用这些说明理解变量含义，不读取说明中已有算法参数。

\clearpage
\section{候选方法、自由度与最终状态}

\begin{longtable}{p{4.0cm} p{2.2cm} p{3.0cm} p{3.0cm}}
\caption{当前方法清单。}
\label{tab:method-index}\\
\toprule
方法 & 典型自由度 & 最终状态 & 主要原因 \\
\midrule
\endfirsthead
\toprule
方法 & 典型自由度 & 最终状态 & 主要原因 \\
\midrule
\endhead
历史派生分支 + Gaussian/exp & 峰形依赖 & 废弃 & 非独立且 GOF 拒绝 \\
第一版独立箱图原型 & 折间变化 & 废弃 & 多指标不一致、选择不稳 \\
A0 raw & 0 & 基线 & 用于同管线前后比较 \\
A1/A2/A3 单通道 & 低 & 负结果 & 两种拆分均恶化 \\
A4 双通道 & 低 & 主方案 & 物理可解释、双拆分稳定 \\
A5 + file 漂移 & 额外线性项 & 不采用 & 增益不跨拆分稳定 \\
challenger \(dt\) only & 3 & 保留对照 & 收益有限 \\
challenger \(r^2\) only & 3 & 保留对照 & 明显改善但非最优 \\
challenger joint \(dt+r^2\) & 6 & 保留对照 & 稳定改善 \\
challenger \(dt+r^2\rightarrow x/y\) & 12 & 冻结候选 & 数字最好，待跨 run 验证 \\
\bottomrule
\end{longtable}

\clearpage
\section{复现命令、环境与输出索引}

\subsection{Python 环境}

本次复核环境为 Python 3.13.12、NumPy 2.4.2、Pandas 3.0.2、SciPy 1.17.1、Matplotlib 3.10.8。
现有 \path{requirements.txt} 未锁定精确版本，这是复现性限制；下一版应生成锁定环境文件。
随机种子在峰形审计中固定为 10972。

\subsection{主要运行入口}

\begin{lstlisting}
python independent_reconstruction.py
python channel_physics_matrix.py
python basic_experiment_matrix.py
python robust_peak_evaluation.py --help
python report/report_assets.py
\end{lstlisting}

不同脚本具有命令行参数；正式复现应以各脚本 \texttt{--help} 与输出 metadata 为准，
避免把本附录的简写命令当作完整配置。

\subsection{代码与证据文件}

\small
\begin{longtable}{p{5.0cm} p{7.2cm}}
\caption{代码和主要证据文件索引。}
\label{tab:file-index}\\
\toprule
路径 & 内容 \\
\midrule
\endfirsthead
\toprule
路径 & 内容 \\
\midrule
\endhead
\path{independent_reconstruction.py} & 第一版基础变量原型，保留失败路径 \\
\path{channel_physics_matrix.py} & A0--A5 分通道物理消融与 A4 主方案 \\
\path{basic_experiment_matrix.py} & challenger 方法矩阵、内层选择和 OOF 预测 \\
\path{robust_peak_evaluation.py} & 固定窗口、多背景、多初值与参数 bootstrap GOF \\
\path{outputs/variable_audit.csv} & 188 字段缺失、无穷、常数审计 \\
\path{outputs/cut_flag_audit.csv} & 14 个 cut 字段 0/1 统计 \\
\path{outputs/README_SUPERSEDED.md} & 旧派生分支结果为何被取代 \\
\path{independent_outputs/summary.json} & 第一版独立原型结果 \\
\path{channel_physics_outputs/summary.csv} & A0--A5 全事件 OOF 指标 \\
\path{channel_physics_outputs/fold_metrics.csv} & A0--A5 逐折结果和训练侧参数 \\
\path{channel_physics_outputs/paired_file_bootstrap_summary.csv} & A4 文件级 bootstrap \\
\path{channel_physics_outputs/robust_peak_corrected_gof.json} & A4 正式峰形 GOF \\
\path{basic_matrix_outputs/summary_metrics.csv} & challenger 方法矩阵 OOF 指标 \\
\path{basic_matrix_outputs/fold_metrics.csv} & challenger 逐折结果 \\
\path{basic_matrix_outputs/inner_selection.csv} & 每个外折的训练侧选择 \\
\path{basic_matrix_outputs/paired_file_bootstrap_summary.csv} & challenger 文件级 bootstrap \\
\path{basic_matrix_outputs/robust_peak_corrected_gof.json} & challenger 正式峰形 GOF \\
\path{report/report_assets.py} & 本报告四张汇总图的生成脚本 \\
\bottomrule
\end{longtable}
\normalsize

\subsection{LaTeX 公式与报告编译规范}

本报告所有行内公式使用 \verb|\(...\)|，独立公式使用 \texttt{equation} 或 \texttt{align}，
重要公式均带 \verb|\label| 并通过 \verb|\cref| 交叉引用；不使用双美元符号公式。
建议用 XeLaTeX 编译两到三次，以生成正确目录和引用。

\clearpage
\section{结果数字的最后核对表}

\begin{table}[H]
\centering
\caption{摘要、正文和结论必须保持一致的冻结数字。}
\label{tab:frozen-numbers}
\begin{tabular}{llcc}
\toprule
方案 & 拆分 & \(\Rsixtyeight\) raw\(\rightarrow\)corr. & \(\Rninety\) raw\(\rightarrow\)corr. \\
\midrule
A4 & file mod 5 & 7.774\%\(\rightarrow\)6.773\% & 14.740\%\(\rightarrow\)13.262\% \\
A4 & 连续块 & 7.893\%\(\rightarrow\)6.779\% & 14.783\%\(\rightarrow\)13.278\% \\
challenger & file mod 5 & 8.019\%\(\rightarrow\)5.833\% & 15.127\%\(\rightarrow\)11.776\% \\
challenger & 连续块 & 8.125\%\(\rightarrow\)5.973\% & 15.029\%\(\rightarrow\)11.826\% \\
\bottomrule
\end{tabular}
\end{table}

\clearpage
\begin{thebibliography}{99}

\bibitem{technote}
PandaX collaboration internal note,
\emph{TechNote for Energy Resolution Improvement Project Using PandaX-4T Calibration Data},
用户提供的变量说明 PDF。

\bibitem{pandax-bb2nu}
PandaX-4T Collaboration,
\emph{Measurement of Double Beta Decay Half-life of \(^{136}\mathrm{Xe}\) with PandaX-4T Detector},
\href{https://arxiv.org/abs/2205.12809}{arXiv:2205.12809}.

\bibitem{pandax-linearity}
PandaX-4T Collaboration,
\emph{Improvement on the Linearity Response of PandaX-4T with new PMT Bases},
\href{https://arxiv.org/abs/2401.00373}{arXiv:2401.00373}.

\bibitem{pandax-mev-search}
PandaX-4T Collaboration,
\emph{Search for MeV-Scale Axionlike Particles and Dark Photons with PandaX-4T},
\href{https://arxiv.org/abs/2409.00773}{arXiv:2409.00773}.

\bibitem{pandax-0nubb}
PandaX-4T Collaboration,
\emph{Searching for Neutrinoless Double-Beta Decay of \(^{136}\mathrm{Xe}\) with PandaX-4T},
\href{https://arxiv.org/abs/2412.13979}{arXiv:2412.13979}.

\end{thebibliography}

\end{document}
